\documentclass{article}

\usepackage[utf8]{inputenc}     % pdfLaTeX
\usepackage[T1]{fontenc}
\usepackage[magyar]{babel}
\usepackage{amsmath}

% Set page size and margins
% Replace `letterpaper' with `a4paper' for UK/EU standard size
\usepackage[a4paper,top=2cm,bottom=2cm,left=3cm,right=3cm,marginparwidth=1.75cm]{geometry}
%for arrows
\usepackage{textcomp}

\title{Magánhangzótörvények}
\author{Lengyel Zoltán Péter\thanks{Csukás Ibolya tanítása alapján}}
\date{Legutóbb frissítve: 2026. Május 04.}

\begin{document}

\maketitle
\tableofcontents
\newpage

\section{Hangrend törvénye}
\begin{itemize}
    \item a magánhangzók harmóniája a szavakba
    \item Magas hangrendű szavak \begin{itemize}
        \item csak magas magánhangzót tartalmaznak 
        \item pl. eper, díszít, keresgéltünk, stb.
    \end{itemize}
    \item Mély hangrendű szavak \begin{itemize}
        \item csak mély magánhangzót tartalmaznak 
        \item pl. alma, táska, autó, stb.
    \end{itemize}
    \item Vegyes hangrendű szavak \begin{itemize}
        \item magas és mély magánhangót is tartalmaznak
        \item pl. dió, hajó, iskola, stb.
    \end{itemize}
\end{itemize}

\section{Az illeszkedés törvénye}
\begin{itemize}
    \item "Milyen hangrendű szóhoz milyen hangrendű toldalék járul"
\end{itemize}

\vspace{5pt}

\begin{itemize}
    \item Magas hangrendű szóhoz magas hangrendű toldalék járul \begin{itemize}
        \item a háromalakú toldalékok esetében a szó utolsó magánhangzója alapján ajakréses vagy ajakkerekítéses
        \item pl. kerthez, erdőhöz, énekeltek, ültök
    \end{itemize}
    \item Mély hangrendű szóhoz mély toldalék járul \begin{itemize}
        \item pl. padban, kapunál, rúgtok
    \end{itemize}
    \item Vegyes hangrendű szavakhoz mély hangrendű toldalék járul \begin{itemize}
        \item pl. iskolában, papírról, rádióhoz
    \end{itemize}
\end{itemize}

\subsection{Toldalékok alakjuk szerint}
\begin{itemize}
    \item egyalakúak pl. -kor, -ig, -ért, -t
    \item kétalakúak pl. -ban/-ben, -ba/-be, -ról/-ről, nál/-nél \begin{itemize}
        \item egy mély és egy magas alak
    \end{itemize}
    \item háromalakúak pl. -on/-en/-ön, -hoz/-hez/-höz, -szor/-szer/-ször, -tok/-tek/-tök \begin{itemize}
        \item egy mély és két magas (egy ajakréses és egy ajakkerekítéses) alak
    \end{itemize}
\end{itemize}

\section{A hiátus törvénye}
\begin{itemize}
    \item Két magánhangzó között gyakran egy j hangot ejtünk \begin{itemize}
        \item pl. dió, diák, fiú, tea
    \end{itemize}
\end{itemize}

\end{document}